\chapter{Statistica applicata: leggere i risultati}
\label{ch:metodo_statistica}

\epigraph{``Statistics is the grammar of science.''}{Karl Pearson,
attribuito}

\lettrine[lines=3]{D}{opo} aver eseguito una pipeline (o un
batch Monte Carlo), il coordinatore vuole sapere: la
soluzione \`e bilanciata? alcune classi sono
sistematicamente svantaggiate? c'\`e una correlazione fra
qualcosa che ho dato in input e la qualit\`a del risultato?

Per rispondere $\pi$Tantum integra una piccola toolbox di
statistica applicata, accessibile dalla pagina
Diagnostica\index{diagnostica}. Le tre famiglie sono:
\emph{distribuzioni e istogrammi},
\emph{correlazioni e regressioni},
\emph{test di ipotesi} (chi-quadro, Kolmogorov-Smirnov).

\section{Distribuzioni: vedere la forma dei dati}

\begin{defbox}
Un \emph{istogramma}\index{istogramma} di una variabile
$X$ \`e una funzione che, suddiviso il range di $X$ in $b$
\emph{bin} di larghezza uguale, conta quanti valori cadono
in ognuno. \`E una stima non parametrica della densit\`a di
$X$.
\end{defbox}

In $\pi$Tantum si possono richiedere istogrammi di:
\begin{itemize}
\item Ore di lezione per docente (mostra se i carichi sono
  distribuiti uniformemente o se ci sono outlier).
\item ``Buchi'' per docente nella settimana (ore vuote fra
  lezioni).
\item Slot per materia per ora del giorno (mostra
  l'effetto della preferenza ``mattutino'').
\item Tempo di esecuzione per cluster nella decomposizione
  spettrale.
\end{itemize}

\subsection{Bin sizing}

\sidenote{La scelta del numero di bin \`e un piccolo
problema in s\'e: troppi pochi $\to$ perdi struttura;
troppi $\to$ rumore.}
$\pi$Tantum usa la regola di Freedman-Diaconis come default:
\[
  h = 2 \cdot \mathrm{IQR}(X) \cdot n^{-1/3}
\]
dove IQR \`e l'\emph{interquartile range} ($Q_3 - Q_1$) e
$n$ il numero di osservazioni. Il numero di bin \`e
$\lceil (\max - \min) / h \rceil$. L'utente pu\`o sempre
forzare un numero specifico.

\subsection{Test di Kolmogorov-Smirnov}

\begin{defbox}
Il \emph{test KS}\index{Kolmogorov-Smirnov!test} confronta
una distribuzione empirica con una distribuzione di
riferimento (o due distribuzioni empiriche fra loro). La
statistica del test \`e:
\[
  D = \sup_x |F_n(x) - F(x)|
\]
ovvero la massima distanza verticale fra la CDF empirica
$F_n$ e la CDF di riferimento $F$.
\end{defbox}

In $\pi$Tantum il test KS \`e usato ad esempio per
confrontare la distribuzione delle ore-per-docente fra
``baseline'' e ``post-perturbazione Monte Carlo''. Un valore
$D$ piccolo (con $p$-value alto) indica che la distribuzione
\`e robusta; un $D$ grande indica che la perturbazione ha
significativamente alterato il bilancio.

\section{Correlazioni e regressioni}

\subsection{Correlazione di Pearson}

\begin{defbox}
La \emph{correlazione di Pearson}\index{Pearson!correlazione}
fra due variabili $X, Y$ \`e:
\[
  r_{XY} = \frac{\mathrm{Cov}(X, Y)}{\sigma_X \sigma_Y}
\]
con $\mathrm{Cov}(X, Y) = E[(X - \bar X)(Y - \bar Y)]$. Vale
fra $-1$ e $+1$: $r > 0$ relazione positiva, $r < 0$
negativa, $r \approx 0$ nessuna relazione lineare.
\end{defbox}

\textbf{Come si legge}: $r$ vicino a $\pm 1$ indica
relazione \emph{lineare} forte; $r$ vicino a $0$ indica
\emph{nessuna relazione lineare}, ma non esclude relazioni
non lineari. Un valore di $r = 0.5$ significa: \emph{circa il
25\% della varianza di $Y$ \`e spiegato linearmente da $X$}
($r^2 = 0.25$).

\begin{esempiobox}
Calcoliamo la correlazione fra ``numero di indisponibilit\`a
HARD del docente'' e ``numero di buchi nel suo orario
finale''. Se $r = +0.6$, significa che pi\`u un docente ha
indisponibilit\`a, pi\`u tende ad avere buchi. Insight per
il coordinatore: i docenti con molte indisponibilit\`a si
prendono comunque l'orario peggiore in termini di buchi;
forse vale la pena rivedere le loro indisponibilit\`a o
accettare di fare strappi sulle preferenze altrui.
\end{esempiobox}

\subsection{Regressione lineare (OLS)}

\begin{defbox}
Una \emph{regressione lineare}
\emph{ordinaria}\index{regressione!OLS} (OLS) modella $Y$
come funzione lineare di una o pi\`u $X$:
\[
  Y = \beta_0 + \beta_1 X_1 + \beta_2 X_2 + \cdots + \beta_p X_p + \epsilon
\]
e stima i $\beta$ minimizzando la somma dei quadrati dei
residui: $\hat{\beta} = (X^\top X)^{-1} X^\top Y$.
\end{defbox}

\textbf{Come si legge}: ogni coefficiente $\beta_i$ dice di
quanto cambia $Y$ se $X_i$ aumenta di $1$, tenendo le altre
$X_j$ fisse. Il coefficiente $\beta_0$ \`e l'\emph{intercetta}.
$\pi$Tantum riporta anche $R^2$ (frazione di varianza
spiegata) e i $p$-value per ciascun $\beta_i$.

In $\pi$Tantum la pagina ``Correlazioni e regressioni''
permette di costruire interattivamente il modello (scegli
le $X$, scegli la $Y$, scegli OLS o Logit) e mostra i
coefficienti con intervalli di confidenza.

\subsection{Regressione logistica (Logit)}

\begin{defbox}
La \emph{regressione logistica}\index{regressione!logistica}
modella la probabilit\`a di un evento binario:
\[
  P(Y = 1 | X) = \frac{1}{1 + e^{-(\beta_0 + \beta_1 X_1 + \cdots)}}
\]
\end{defbox}

Si usa quando la $Y$ \`e binaria. Esempio in $\pi$Tantum:
``probabilit\`a che una classe abbia $\geq 3$ buchi
settimanali in funzione del suo numero di indirizzi serviti
e del numero di docenti coinvolti''.

\section{Chi-quadro: indipendenza}

\begin{defbox}
Il \emph{test chi-quadro
$\chi^2$}\index{chi-quadro!test} verifica se due variabili
categoriali sono indipendenti. Si costruisce la tabella di
contingenza, si calcola
\[
  \chi^2 = \sum_{i,j} \frac{(O_{ij} - E_{ij})^2}{E_{ij}}
\]
dove $O_{ij}$ \`e la frequenza osservata e $E_{ij}$ quella
attesa sotto ipotesi di indipendenza. Si confronta con la
distribuzione $\chi^2$ con $(r-1)(c-1)$ gradi di libert\`a.
\end{defbox}

Tipico uso: ``il giorno della settimana e la materia sono
indipendenti?'' Se p-value basso, la materia \emph{influisce}
su quale giorno si tiene (es. matematica concentrata al
mattino). \`E un'indicazione che il sistema sta rispettando
la preferenza ``mattutino''.

\section{Visualizzazione interattiva}

Tutti i grafici in $\pi$Tantum sono prodotti con
ECharts~6\index{ECharts}, lazy-loaded sul frontend. Hanno
sempre i bottoni:
\begin{itemize}
\item \textbf{Export PNG} (con pixel-ratio 2 e sfondo
  bianco).
\item \textbf{Clear graph}: per fare spazio in finestra.
\item \textbf{Restore zoom}: per resettare lo zoom dopo
  esplorazione.
\end{itemize}

I parametri (numero di bin, scelta delle $X$, etc.) sono
modificabili dall'UI senza dover rilanciare il run; le
visualizzazioni sono ricalcolate al volo dal browser.

\section{Connessioni}

\begin{connessionebox}
\begin{itemize}
\item \textbf{Monte Carlo} (cap.~\ref{ch:metodo_montecarlo}):
  produce campioni che si analizzano con distribuzioni
  e KS-test.
\item \textbf{Analisi del grafo}
  (cap.~\ref{ch:metodo_grafo}): la modularit\`a e la
  betweenness sono input naturali per le regressioni
  (es. ``il tempo di esecuzione del solver dipende dalla
  modularit\`a del grafo?'').
\item \textbf{Workflow}: dopo ogni esecuzione il sistema
  popola automaticamente le tabelle delle statistiche
  (lezioni, buchi, ore-isolate, etc.); il coordinatore le
  pu\`o esplorare immediatamente.
\end{itemize}
\end{connessionebox}

\begin{biblioparvabox}
\begin{itemize}
\item \cite{wasserman2004}: \emph{All of Statistics},
  manuale conciso di statistica matematica accessibile
  agli informatici.
\item \cite{james2013}: \emph{An Introduction to
  Statistical Learning}, gratuito online, ottimo per
  regressioni e modelli predittivi.
\end{itemize}
\end{biblioparvabox}
